\chapter{Introduction}

\section{Introduction}

OrchardVision is a computer vision system for long-term fruit tracking and counting in mango orchards. It implements identity learning, tracking pipelines, and a full-stack web interface for experiment management and video processing. A user uploads orchard video through the web interface; the backend extracts frames, detects fruit, tracks each fruit through time, learns a re-identification (ReID) embedding for it, consolidates the observations into persistent identities and finally computes analytics such as IDF1, HOTA and the number of identity switches.

\add{Mango is one of the most widely grown tropical fruit crops, and the timing of harvest, labour planning, packing-house capacity and market commitments all depend on how many fruit an orchard will deliver. Growers traditionally estimate this by hand-counting a small sample of trees and extrapolating to the whole block. The process is slow, subjective and statistically weak, because fruit load varies strongly from tree to tree and from row to row.}

\add{Video captured by a handheld camera, a tractor-mounted camera or an unmanned aerial vehicle offers a much denser sample of the orchard. Turning that video into a reliable count, however, is not simply a matter of running an object detector on every frame. Each fruit is visible in dozens or hundreds of consecutive frames, and a detector will report it every time. The central problem is therefore one of \emph{identity}: deciding whether a detection in the current frame is a fruit that has already been counted or a new one.}

\section{Overview of Orchard Fruit Counting}

\add{Automatic fruit counting systems typically combine three ingredients. A detector finds fruit in each frame; a tracker links detections of the same fruit across consecutive frames into short trajectories; and a counting rule converts trajectories into a number, for example by counting the number of unique trajectories or by counting crossings of a virtual line. This design works acceptably when fruit are always visible and the camera moves smoothly in one direction.}

\add{Real orchards violate both assumptions. Foliage and branches hide fruit for tens of frames at a time, which breaks trajectories and inflates the count when the fruit reappears. In long-term monitoring the same tree is recorded on several days, or the camera returns to a row after a turn, so identities must survive gaps of seconds, minutes or days. OrchardVision addresses this by adding an \emph{identity learning} layer on top of tracking: every identity is represented by a prototype embedding that is updated as new observations arrive, and tracks that end and later reappear can be linked back to the same identity through a ReID step.}

The system is organised as a pipeline (frame extraction, detection, tracking, ReID embedding, identity consolidation, analytics), a REST API that exposes every intermediate and final result, a relational data model that persists them, and a web client with an experiment harness for comparing configurations.

\section{Challenges in Orchard Video}

\add{The following factors make fruit tracking in orchards considerably harder than the pedestrian and vehicle tracking problems on which most multi-object tracking research has been developed:}

\begin{itemize}
  \item \add{\textbf{Occlusion.} Leaves, branches and other fruit frequently hide a fruit partially or completely, so tracks fragment.}
  \item \add{\textbf{Similar appearance.} Fruit on the same tree share colour, shape and texture, which weakens appearance cues that ReID relies on.}
  \item \add{\textbf{Camera motion.} Handheld or vehicle-mounted cameras move constantly, so image-space motion models (such as constant-velocity Kalman filters) are disturbed by ego-motion rather than by object motion.}
  \item \add{\textbf{Illumination and scale changes.} Sunlight, shadows and varying distance to the canopy change the apparent size and colour of fruit.}
  \item \add{\textbf{Long gaps.} In long-term monitoring an identity may disappear for far longer than any tracker buffer, so matching must rely on learned appearance rather than short-term motion.}
  \item \add{\textbf{Dense small objects.} A single frame may contain tens of small, closely spaced fruit, increasing the risk of identity switches.}
\end{itemize}

\section{Motivation}

\add{Three observations motivate the project. First, yield estimates are needed early and repeatedly during the season, and manual sampling cannot keep up with this demand at the scale of commercial orchards. Second, video is cheap to collect and is already captured during routine orchard operations, so a counting method that runs on ordinary video is easy to adopt. Third, counting errors from video are dominated by identity errors rather than detection errors, which means that progress requires explicit support for identity: storing what each fruit looks like, remembering it after it disappears and measuring how well this works.}

\add{A further motivation is engineering. Research prototypes in this area are often a script that prints a number. A reusable system needs persistent results, a stable API, reproducible experiments and a user interface, so that agronomists and researchers can examine and trust the output. OrchardVision therefore treats the data model, the API and the experiment harness as first-class parts of the contribution.}

\section{Project Statement}

\add{To design and implement an end-to-end, web-based system that ingests orchard video, detects and tracks individual mango fruit, maintains a persistent learned identity for each fruit across occlusions and re-appearances, and reports tracking quality with standard metrics, while making every stage configurable and every result inspectable and comparable through a reproducible experiment harness.}

\section{Objectives}

The objectives of the project are:

\begin{enumerate}
  \item To implement a processing pipeline that orchestrates frame extraction, detection, tracking, ReID embedding, identity consolidation and analytics for an uploaded video.
  \item To support multiple tracker back-ends (mock, ByteTrack and TAPNext) behind a common interface so that they can be compared on the same video.
  \item To learn and persist an identity prototype (embedding vector, quality and observation count) for every fruit and to record ReID events when an identity is lost and later reappears.
  \item To estimate camera motion and record whether compensation was applied.
  \item To compute IDF1, HOTA and ID-switch metrics and per-frame timing telemetry.
  \item To expose all results through a documented REST API and a React~/~TypeScript web client with shared type definitions.
  \item To provide an experiment harness that runs and stores tracker evaluations and exports them as JSON or CSV.
  \item \add{To make the system deployable on a developer laptop (SQLite) and in production (PostgreSQL, object storage, background workers).}
\end{enumerate}

\section{Scope of the Project}

The scope of the work covers the backend service, the data model, the processing pipeline, the REST API, the shared TypeScript types, the deployment configuration and the experiment-management interface described in the following chapters. Uploads are limited to 2\,GB per video by default (configurable through the \code{MAX_UPLOAD_BYTES} environment variable).

\add{The following are explicitly outside the scope of this report: training new fruit detectors from scratch, fruit-size or ripeness estimation, three-dimensional orchard mapping, and multi-camera fusion. The evaluation reported in Chapter~8 is limited to the videos and tracker configurations listed there.}

\section{Organisation of the Report}

\add{Chapter~2 reviews the literature on fruit detection, multi-object tracking, re-identification and tracking metrics. Chapter~3 presents the requirements and the system architecture. Chapter~4 describes the data model. Chapter~5 explains the processing pipeline and the identity-learning design. Chapter~6 documents the REST API, error handling and shared types. Chapter~7 covers deployment and testing. Chapter~8 presents the evaluation set-up and results, and Chapter~9 concludes the report and outlines future work. References and appendices follow.}
